\documentclass[10pt,a4paper]{amsart}
\usepackage{amsmath,amssymb,amsthm,graphicx}
\usepackage{verbatim}
\newenvironment{reference}{\comment}{\endcomment}
%\newenvironment{reference}{}{}
\newenvironment{slogan}{\comment}{\endcomment}
\newenvironment{history}{\comment}{\endcomment}

% For commutative diagrams we use Xy-pic
\usepackage[all]{xy}

% We use 2cell for 2-commutative diagrams.
\xyoption{2cell}
% Theorem environments.
%
\theoremstyle{plain}
\newtheorem{theorem}[subsection]{Theorem}
\newtheorem{proposition}[subsection]{Proposition}
\newtheorem{lemma}[subsection]{Lemma}

\theoremstyle{definition}
\newtheorem{definition}[subsection]{Definition}
\newtheorem{example}[subsection]{Example}
\newtheorem{exercise}[subsection]{Exercise}
\newtheorem{situation}[subsection]{Situation}

\theoremstyle{remark}
\newtheorem{remark}[subsection]{Remark}
\newtheorem{remarks}[subsection]{Remarks}

\numberwithin{equation}{subsection}

% Macros
%
\def\lim{\mathop{\mathrm{lim}}\nolimits}
\def\colim{\mathop{\mathrm{colim}}\nolimits}
\def\Spec{\mathop{\mathrm{Spec}}}
\def\Hom{\mathop{\mathrm{Hom}}\nolimits}
\def\Ext{\mathop{\mathrm{Ext}}\nolimits}
\def\SheafHom{\mathop{\mathcal{H}\!\mathit{om}}\nolimits}
\def\SheafExt{\mathop{\mathcal{E}\!\mathit{xt}}\nolimits}
\def\Sch{\mathit{Sch}}
\def\Mor{\mathop{\mathrm{Mor}}\nolimits}
\def\Ob{\mathop{\mathrm{Ob}}\nolimits}
\def\Sh{\mathop{\mathit{Sh}}\nolimits}
\def\NL{\mathop{N\!L}\nolimits}
\def\CH{\mathop{\mathrm{CH}}\nolimits}
\def\proetale{{pro\text{-}\acute{e}tale}}
\def\etale{{\acute{e}tale}}
\def\QCoh{\mathit{QCoh}}
\def\Ker{\mathop{\mathrm{Ker}}}
\def\Im{\mathop{\mathrm{Im}}}
\def\Coker{\mathop{\mathrm{Coker}}}
\def\Coim{\mathop{\mathrm{Coim}}}

% Boxtimes
%
\DeclareMathSymbol{\boxtimes}{\mathbin}{AMSa}{"02}

%
% Macros for moduli stacks/spaces
%
\def\QCohstack{\mathcal{QC}\!\mathit{oh}}
\def\Cohstack{\mathcal{C}\!\mathit{oh}}
\def\Spacesstack{\mathcal{S}\!\mathit{paces}}
\def\Quotfunctor{\mathrm{Quot}}
\def\Hilbfunctor{\mathrm{Hilb}}
\def\Curvesstack{\mathcal{C}\!\mathit{urves}}
\def\Polarizedstack{\mathcal{P}\!\mathit{olarized}}
\def\Complexesstack{\mathcal{C}\!\mathit{omplexes}}
% \Pic is the operator that assigns to X its picard group, usage \Pic(X)
% \Picardstack_{X/B} denotes the Picard stack of X over B
% \Picardfunctor_{X/B} denotes the Picard functor of X over B
\def\Pic{\mathop{\mathrm{Pic}}\nolimits}
\def\Picardstack{\mathcal{P}\!\mathit{ic}}
\def\Picardfunctor{\mathrm{Pic}}
\def\Deformationcategory{\mathcal{D}\!\mathit{ef}}
\usepackage{cleveref}
\crefname{originallemma}{original lemma}{original lemmas}
\Crefname{originallemma}{Original lemma}{Original lemmas}
\makeatletter
\newlabel{algebra-lemma-derivative-zero-pth-power}{{031W}{source}{}{sourcepointer-1}{}}
\newlabel{algebra-lemma-derivative-zero-pth-power@cref}{{[originallemma][0][]031W}{[1][0][]source}{}{}{}}
\makeatother
\title{A quadratic element of the square of an ideal of large colength is a square up to a factor}
\author{The Stacks Project Authors}
\date{stacks-project:resolve.tex:local-snapshot-20261012; GNU FDL 1.2 or later}
\begin{document}
\maketitle
\noindent\textit{Editorial note.} All delivered mathematical text is the author's own, copied from the pinned Stacks Project chapter file resolve.tex; no statement or proof step was rewritten, reordered or removed.\par\medskip
\noindent\textit{Editorial note.} The counted proof is the author's complete argument: the reduction in the first case through the partial derivatives and the comparison of degrees inside the square of the ideal, the characteristic two analysis using derivations of the field to produce polynomials of the shape constant plus square in the ideal, and the final exclusion through the conormal module and a dimension count with an explicit coordinate change.\par\medskip
\noindent\textit{Editorial note.} The label lemma-issquare names the counted statement, and the proof environment immediately following it in the source is delivered verbatim as its proof.\par\medskip
\noindent\textit{Editorial note.} The counted statement is carried by the author's own lemma environment in the chapter source.\par\medskip
\noindent\textit{Editorial note.} Pre-open gate: python operations/agent-authoring/tools/precard\_gate.py --source resolve.tex --label lemma-issquare --start 3315 --end 3398 --local-id 28070 returned RESULT=OK (L28070 free, occupancy FREE, public interval clear). Pre-assembly tag self-check: canonical tag 0BGD, not present in reference/public-items.jsonl (neither by tag\_url nor by claim\_key suffix).\par\medskip
\noindent\textit{Editorial note.} The statement and proof are copied from the pinned chapter file; the cited inputs are kept as external pointers with their canonical Stacks tags.\par\medskip
\noindent\textit{Allowed standard prerequisites (editorial).}\begin{itemize}
\item Leibniz rule for the partial derivatives acting on the square of an ideal. Standard calculus of partial derivatives in a polynomial ring: the derivative of an element of the square of an ideal lies in the ideal. Supplies the fact that the two partial derivatives of the given quadratic element lie in the ideal, which drives the first case distinction.
\item Extension of a derivation of the field with nontrivial value. Algebra, the lemma that for an element which is not a p-th power there is a derivation taking a nonzero value on it, extended to the polynomial ring by sending the variables to zero. Supplies the derivation used in characteristic two to produce elements of the ideal of a controlled shape.
\item Freeness of the conormal module of the auxiliary ideal. Standard description of the conormal module of an ideal generated by a regular sequence of two elements in a two dimensional situation. Supplies the conclusion that the inclusion of the auxiliary ideal in the given one is strict, which produces an element with a nonzero mixed coefficient.
\end{itemize}\medskip
\ifdefined\hypertarget\hypertarget{sourcepointer-1}{}\fi
\noindent\textit{Source reference.} Original lemma 031W: sources/stacks-chow/tags.txt, source lines 3999--3999. Prerequisite pointer invoked at proof line(s) 3363; it names a standard fact of a sibling chapter that the counted argument uses, and every step of the counted proof built on it is retained verbatim in the delivered body.\par\medskip
\begin{lemma}
\label{lemma-issquare}
Let $\kappa$ be a field. Let $I \subset \kappa[x, y]$ be an ideal. Let
$$
a + b x + c y + d x^2 + exy + f y^2 \in I^2
$$
for some $a, b, c, d, e, f \in k$ not all zero. If the colength
of $I$ in $\kappa[x, y]$ is $> 1$, then
$a + b x + c y + d x^2 + exy + f y^2 = j(g + hx + iy)^2$
for some $j, g, h, i \in \kappa$.
\end{lemma}

\begin{proof}
Consider the partial derivatives $b + 2dx + ey$ and
$c + ex + 2fy$. By the Leibniz rules these are contained in $I$.
If one of these is nonzero, then after a linear change of coordinates,
i.e., of the form $x \mapsto \alpha + \beta x + \gamma y$ and
$y \mapsto \delta + \epsilon x + \zeta y$, we may assume
that $x \in I$. Then we see that $I = (x)$ or $I = (x, F)$ with
$F$ a monic polynomial of degree $\geq 2$ in $y$.
In the first case the statement is clear. In the second case
observe that we can write any element in $I^2$ in the form
$$
A(x, y) x^2 + B(y) x F + C(y) F^2
$$
for some $A(x, y) \in \kappa[x, y]$ and $B, C \in \kappa[y]$.
Thus
$$
a + b x + c y + d x^2 + exy + f y^2 = A(x, y) x^2 + B(y) x F + C(y) F^2
$$
and by degree reasons we see that $B = C = 0$ and $A$ is a constant.

\medskip\noindent
To finish the proof we need to deal with the case that both
partial derivatives are zero. This can only happen in characteristic $2$
and then we get
$$
a + d x^2 + f y^2 \in I^2
$$
We may assume $f$ is nonzero (if not, then switch the roles of $x$ and $y$).
After dividing by $f$ we obtain the case where the characteristic of
$\kappa$ is $2$ and
$$
a + d x^2 + y^2 \in I^2
$$
If $a$ and $d$ are squares in $\kappa$, then we are done. If not,
then there exists a derivation $\theta : \kappa \to \kappa$ with
$\theta(a) \not = 0$ or $\theta(d) \not = 0$, see
Algebra, Lemma \ref{algebra-lemma-derivative-zero-pth-power}.
We can extend this to a derivation of $\kappa[x, y]$ by setting
$\theta(x) = \theta(y) = 0$. Then we find that
$$
\theta(a) + \theta(d) x^2 \in I
$$
The case $\theta(d) = 0$ is absurd. Thus we may assume
that $\alpha + x^2 \in I$ for some $\alpha \in \kappa$.
Combining with the above we find that $a + \alpha d + y^2 \in I$.
Hence
$$
J = (\alpha + x^2, a + \alpha d + y^2) \subset I
$$
with codimension at most $2$. Observe that
$J/J^2$ is free over $\kappa[x, y]/J$ with basis
$\alpha + x^2$ and $a + \alpha d + y^2$.
Thus $a + d x^2 + y^2 =
1 \cdot (a + \alpha d + y^2) + d \cdot (\alpha + x^2) \in I^2$
implies that the inclusion $J \subset I$ is strict.
Thus we find a nonzero element of the form $g + hx + iy + jxy$ in $I$.
If $j = 0$, then $I$ contains a linear form and we can
conclude as in the first paragraph. Thus $j \not = 0$
and $\dim_\kappa(I/J) = 1$ (otherwise we could find
an element as above in $I$ with $j = 0$).
We conclude that $I$ has the form
$(\alpha + x^2, \beta + y^2, g + hx + iy + jxy)$
with $j \not = 0$ and has colength $3$.
In this case $a + dx^2 + y^2 \in I^2$ is impossible.
This can be shown by a direct computation, but we prefer to argue
as follows. Namely, to prove this statement we may assume that
$\kappa$ is algebraically closed. Then we can do a coordinate
change $x \mapsto \sqrt{\alpha} + x$ and $y \mapsto \sqrt{\beta} + y$
and assume that $I = (x^2, y^2, g' + h'x + i'y + jxy)$ with the same $j$.
Then $g' = h' = i' = 0$ otherwise the colength of $I$ is not $3$.
Thus we get $I = (x^2, y^2, xy)$ and the result is clear.
\end{proof}

\end{document}
